% Generated draft: requires mathematical and source review.
\documentclass[11pt]{article}
\usepackage[margin=25mm]{geometry}
\usepackage{fontspec}
\setmainfont{Latin Modern Roman}
\usepackage{amsmath,amssymb,mathtools,unicode-math}
\setmathfont{Latin Modern Math}
\usepackage{enumitem,needspace,xurl,hyperref}
\setlength{\parindent}{0pt}
\setlength{\parskip}{5pt}
\setlength{\emergencystretch}{3em}
\newcommand{\N}{\mathbb{N}}
\newcommand{\Z}{\mathbb{Z}}
\newcommand{\Q}{\mathbb{Q}}
\newcommand{\R}{\mathbb{R}}
\newcommand{\C}{\mathbb{C}}
\newcommand{\F}{\mathbb{F}}
\newcommand{\A}{\mathbb{A}}
\newcommand{\PP}{\mathbb{P}}
\newcommand{\E}{\mathbb{E}}
\newcommand{\eps}{\varepsilon}
\newcommand{\dd}{\,\mathrm{d}}
\begin{document}
\section*{UnsolvedMath ID: 30003903}
\textbf{Source problem code:} catalogue entry
\subsection{Definitions and mathematical statement}
The supplied record does not specify the ground field or the convention for cohomological dimension. The following makes explicit the usual complex-geometric interpretation of ``singular K3 surface'' and gives a conditional formulation using ordinary integral group cohomological dimension; these conventions require confirmation against the source.

A complex K3 surface is a smooth projective surface \(X/\C\) with trivial canonical bundle and \(H^1(X,\mathcal O_X)=0\). In this context, \emph{singular} means that its Picard number is \(\rho(X)=20\); it does not mean that \(X\) has singular points. Let \(\operatorname{NS}(X)\) be its N\'eron--Severi lattice, equipped with the intersection pairing, and let
\[
T(X)=\operatorname{NS}(X)^{\perp}\subset H^2(X,\Z)
\]
be its transcendental lattice. For such a surface, \(T(X)\) is a positive definite lattice of rank two. Interpret the positive discriminant as
\[
d(X)=\det\bigl((t_i\cdot t_j)_{i,j=1}^{2}\bigr)
       =\bigl|\operatorname{disc}\operatorname{NS}(X)\bigr|,
\]
where \((t_1,t_2)\) is an integral basis of \(T(X)\).

Let \(G=\operatorname{Aut}_{\C}(X)\) be the group of biregular automorphisms, regarded as a discrete group. Its ordinary integral cohomological dimension is
\[
\operatorname{cd}_{\Z}(G)
=\operatorname{pd}_{\Z[G]}\Z,
\]
where \(\Z\) is the trivial \(\Z[G]\)-module and the projective dimension is allowed to be infinite.

\textbf{Conditional formulation of the supplied question.} Under these conventions, is it true that
\[
\forall X/\C,\qquad
\bigl(X\text{ is a singular K3 surface and }d(X)\in\{7,8\}\bigr)
\Longrightarrow
\operatorname{cd}_{\Z}\!\left(\operatorname{Aut}_{\C}(X)\right)=2?
\]
The original record states the equality declaratively, while its cleaned statement asks whether it holds for every such surface. It is treated here as a proposed equality requiring proof.

\textbf{Unresolved qualification.} The source record says only ``cohomological dimension.'' Ordinary integral cohomological dimension, cohomological dimension with other coefficients, and virtual cohomological dimension are different invariants. For a nonzero commutative ring \(R\) with identity, the coefficient-dependent invariant is \(\operatorname{cd}_{R}(G)=\operatorname{pd}_{R[G]}R\). If virtual cohomological dimension is intended, one instead uses the integral cohomological dimension of a torsion-free subgroup of finite index, subject to existence of such a subgroup. The record does not identify the intended invariant, coefficient ring, or any restriction on the automorphism group. Consequently, the displayed integral formulation is an explicit interpretation, not a verified repair of the original question.
\subsection{Short English statement}
The proposed equality concerns every singular K3 surface of positive discriminant \(7\) or \(8\). Under the complex interpretation, these are smooth K3 surfaces of Picard number \(20\). The question asks whether their automorphism groups have cohomological dimension \(2\). The ground field and the meaning of cohomological dimension require source confirmation, especially the distinction between ordinary and virtual cohomological dimension.
\subsection{Sources}
\begin{itemize}
\item Local UnsolvedMath registry, permanent ID \(30003903\), source problem code \texttt{catalogue entry}, title ``Cohomological Dimension of Automorphism Groups of Singular K3 Surfaces,'' in the supplied Oberwolfach Reports Open Problems collection. No individual proposer is supplied; the proposed year is recorded as \(2018\).
\item Supplied report citation: \emph{Classical Algebraic Geometry (2019)}, Oberwolfach Reports, pp.\ 1983--2031, DOI \texttt{10.4171/owr/2018/33}. Supplied link: \url{https://doi.org/10.4171/owr/2018/33}. The record describes these reports as documenting mathematical workshops and their open-problem discussions. The citation's \(2019\) label and the record's \(2018\) dating are retained without reconciliation.
\item Supplied literature source labelled ``Ujikawa'': \url{https://arxiv.org/abs/1307.4197}.
\item Supplied literature source labelled ``Shimada'': \url{https://arxiv.org/abs/1304.7427}.
\item Supplied literature source labelled ``OWR report'': \url{https://ems.press/content/serial-article-files/46758?nt=1}.
\item Supplied literature assessment, checked \texttt{2026-08-21}: the curation reports that Ujikawa's generators and Bertin--Lecacheux's elliptic-fibration classifications provide substantial data, that no proof of the stated equality was found, and that this special case remains open. This is supplied dated curation derived from the OWR material, not an independently verified literature assessment. No specific Bertin--Lecacheux reference is supplied. The record expressly requests independent verification of time-sensitive status and cited claims.
\item The registry labels the problem \texttt{open} and the statement \texttt{corrected\_verified}, explaining that notation, typography, and omitted definitions were repaired while preserving mathematical content. These are supplied registry labels, not mathematical verification performed for this draft.
\end{itemize}
\end{document}
